\noindent\textbf{从磁偶极矩密度到束缚电流。}
磁化强度 $\vect M$ 是单位体积的磁偶极矩。令
$R=|\vect r-\vect r'|$，体元的偶极矢势叠加为
\[
 \vect A_M(\vect r)=\frac{\mu_0}{4\pi}\int_V
 \vect M(\vect r')\times\nabla'\frac1R\dd V'.
\]
由 $\nabla'\times(\vect M/R)
=\nabla'(1/R)\times\vect M+(\nabla'\times\vect M)/R$，
并把旋度体积分转为边界积分，得到
\[
 \vect A_M=\frac{\mu_0}{4\pi}\left[
 \int_V\frac{\nabla'\times\vect M}{R}\dd V'
 +\oint_{\partial V}\frac{\vect M\times\vect n}{R}\dd S'\right].
\]
与电流的矢势比较，便识别出
\[
 \boxed{\vect j_M=\nabla\times\vect M,\qquad
 \vect K_M=\vect M\times\vect n}.
\]
若界面两侧都有磁化，$\vect n$ 从介质1指向介质2，则净面流为
$(\vect M_1-\vect M_2)\times\vect n$。
均匀磁化体内部 $\vect j_M=0$，并不意味着没有磁场，边界面流仍然存在。

\begin{qualityfigure}
\begin{tikzpicture}[>=Latex,font=\small,scale=.9]
 \fill[PaleBlue](-1.6,-1.1)rectangle(1.3,1.35);
 \draw[thick](1.3,-1.1)--(1.3,1.35);
 \foreach\x in{-.8,.2}{\draw[->,DeepBlue,thick](\x,-.65)--(\x,.7);}
 \node[left]at(-.8,.1){$\vect M$};
 \draw[->](1.3,.9)--(2.45,.9)node[right]{$\vect n=\vect e_x$};
 \node[draw,circle,inner sep=1pt]at(1.3,-.35){$\times$};
 \node[right]at(1.7,-.35){$\vect K_M=M\vect e_y$（入纸面）};
 \draw[->](-1.8,-1.65)--(-.65,-1.65)node[right]{$x$};
 \draw[->](-1.8,-1.65)--(-1.8,-.6)node[above]{$z$};
\end{tikzpicture}
\figcap{右手坐标的 $xz$ 截面：$+y$ 指向纸内，因此 $\vect e_z\times\vect e_x=\vect e_y$ 的面流画作叉号。\label{complete:em:bound-current}}
\end{qualityfigure}

图~\ref{complete:em:bound-current} 把 $\vect M$、外法线和面流三个方向同时固定，
可用于检查下一例磁化球的面流符号。定义
$\vect H=\vect B/\mu_0-\vect M$ 后，静磁 Maxwell 方程成为
\[
 \nabla\times\vect H=\vect j_f,\qquad
 \nabla\cdot\vect H=-\nabla\cdot\vect M.
\]
在无自由电流的单连通区域才能全局写 $\vect H=-\nabla\psi_M$；于是
\[
 \rho_M=-\nabla\cdot\vect M,\qquad
 \sigma_M=(\vect M_1-\vect M_2)\cdot\vect n,
\]
\[
 \nabla^2\psi_M=-\rho_M,\qquad
 \psi_M(\vect r)=\frac1{4\pi}\left[
 \int\frac{\rho_M}{R}\dd V'+\oint\frac{\sigma_M}{R}\dd S'\right].
\]
这些磁荷只是计算 $\vect H$ 的等效源；实际 $\vect B=\mu_0(\vect H+\vect M)$
仍满足 $\nabla\cdot\vect B=0$，不能把等效磁荷理解为真实磁单极。
若有自由电流，可先求其真空形式的场
\[
 \vect H_f=\frac1{4\pi}\int
 \frac{\vect j_f(\vect r')\times(\vect r-\vect r')}{R^3}\dd V',
 \qquad \vect H=\vect H_f-\nabla\psi_M,
\]
再联立磁化关系与边界条件确定 $\vect M$。

对准静态、无磁滞且可逆的介质，固定材料参数时增量能量为
$\dd w=\vect H\cdot\dd\vect B$。
若本构关系是通过原点的线性各向同性关系
$\vect M=\chi_m\vect H$、$\vect B=\mu\vect H$，从零场积分得到
\[
 w=\int_0^{\vect B}\vect H(\vect B')\cdot\dd\vect B'
 =\frac12\vect B\cdot\vect H.
\]
有永久磁化、磁滞或色散时，不能不加条件地套用这一简单能量式。

